\section{Mapa y definiciones: no confundir las tres capas Base, Instruction y Preference}

La sección anterior explicó por qué surgió la alineación; esta sección explica de qué se habla. La comunidad abierta suele usar como sinónimos ``fine-tuned'', ``chat'', ``instruct'', ``aligned'' y ``RLHF model''; el curso, en cambio, exige distinguir entre datos de entrenamiento, función de pérdida e interfaz del modelo. Solo después de separar las capas tienen sentido las comparaciones de modelos y las evaluaciones posteriores.

\subsection{El atlas del curso y los cuatro capítulos}

Esta sección revisa primero la colección de modelos y el atlas de capítulos que proporciona el ponente. Los íconos de modelos no son una genealogía completa, sino una pista reproducible: el lector puede encontrar muchos de los modelos de la clase en la Hugging Face collection. Los cuatro capítulos se desarrollan en torno a instruction, evaluation, RLHF y modern ecosystem, lo que indica que, además de los algoritmos, existe una infraestructura de datos y de medición.

\lecturefigure{slide-017.jpg}{Atlas del curso: colección de modelos, recursos de referencia y línea de tiempo de lanzamientos}{official deck, p.~17}

\lecturefigure{slide-019.jpg}{Enfoque del primer capítulo: instruction models}{official deck, p.~19}

\lecturefigure{slide-020.jpg}{El segundo capítulo agrega: expectations y evaluations}{official deck, p.~20}

\lecturefigure{slide-021.jpg}{El tercer capítulo agrega: RLHF y preference optimization}{official deck, p.~21}

\lecturefigure{slide-022.jpg}{El cuarto capítulo agrega: modern ecosystem}{official deck, p.~22}

\lecturefigure{slide-023.jpg}{Mapa completo de capítulos: instruction, evals, RLHF, ecosystem}{official deck, p.~23}

\begin{knowledgebox}{Lectura de la figura: por qué evaluation ocupa un capítulo propio}
Sin una evaluación rápida, un equipo abierto no puede saber si un cambio en los datos o en la optimización merece publicarse; sin una señal de preferencia humana a largo plazo, la evaluación automática rápida termina siendo manipulada. Evaluation no es un paso accesorio posterior al entrenamiento, sino el canal de retroalimentación que determina si el equipo puede iterar, qué modelos se vuelven visibles y qué datos se siguen recolectando.
\end{knowledgebox}

\subsection{Base model y aligned model}

Ahora dividimos los modelos en dos capas. El base model aprende principalmente una distribución general mediante next-token prediction a gran escala; el modelo que pasa por alignment, fine-tuning o preference training continúa entrenándose sobre él para que sus salidas se adapten al formato de instrucciones, a los roles de diálogo y a las preferencias. Este segundo modelo puede cambiar la experiencia de interacción con relativamente poco cómputo, pero su techo de capacidades sigue fuertemente limitado por el primero.

\lecturefigure{slide-024.jpg}{Los base models son la base de capacidades del ecosistema abierto de alineación}{official deck, p.~24}

\lecturefigure{slide-025.jpg}{Los aligned, fine-tuned y preference-trained models son la capa de adaptación de la conducta}{official deck, p.~25}

\begin{importantbox}{Diferencia de probabilidad condicional entre base y aligned}
El base model aprende $p_{\theta_0}(x_t\mid x_{<t})$; la etapa de alineación produce los parámetros $\theta$, que reordenan la probabilidad de las respuestas bajo la distribución de interacción objetivo. Conceptualmente, puede escribirse
\[
p_{\theta}(y\mid x)\propto p_{\theta_0}(y\mid x)\exp\!\left(\frac{r(x,y)}{\beta}\right),
\]
donde $r(x,y)$ es la preferencia o recompensa asignada a la respuesta y $\beta$ controla cuánto se aleja el modelo de la base/reference distribution. La fórmula expresa el objetivo de «conservar las capacidades y, al mismo tiempo, favorecer las respuestas preferidas»; no significa que todos los métodos lo implementen explícitamente de esta forma.
\end{importantbox}

\teachervoice{00:08:48--00:09:38: el ponente reconoce por iniciativa propia que la línea de tiempo omite muchas contribuciones académicas y de infraestructura, y subraya que sin base models como LLaMA/Llama 2 no existiría el ecosistema abierto de alineación posterior. Un alignment model no es un producto que exista independientemente de su base.}

\subsection{Glosario de alineación}

Esta sección descompone los términos frecuentes según «datos—objetivo—salida». Instruction fine-tuning y supervised fine-tuning suelen usar la misma autoregressive loss, pero el primero pone el énfasis en la interfaz de instrucciones; RLHF usa preferencias humanas para aprender una recompensa y optimizar la policy; DPO optimiza directamente el log-ratio policy/reference a partir de pares de preferencia. Alignment, por su parte, es un concepto de objetivo más amplio.

\lecturefigure{slide-026.jpg}{Definiciones de clase de instruction fine-tuning, SFT, RLHF, DPO y alignment}{official deck, p.~26}

\begin{knowledgebox}{Términos clave: nombre, datos y mecanismo}
\begin{tabularx}{\textwidth}{L{0.18\textwidth} L{0.24\textwidth} X}
\toprule
Término & Datos principales & Mecanismo central \\
\midrule
IFT & instruction-response pairs & Usa la LM loss para aprender el formato de instrucciones, el system prompt, los roles de varios turnos y el estilo objetivo. \\
SFT & Cualquier demostración etiquetada & Usa una pérdida supervisada para imitar la salida objetivo; no necesariamente se trata de instrucciones en lenguaje natural. \\
RLHF & preference comparisons + prompts & Aprende un reward model y luego usa métodos como PPO para maximizar la recompensa y limitar la policy drift. \\
DPO & chosen/rejected response pairs & Optimiza directamente el log-ratio de preferencia, sin ejecutar explícitamente un reward-model rollout loop. \\
Alignment & Estándares de conducta que esperan los usuarios o la organización & Puede lograrse combinando diversos datos, pérdidas, reglas y mecanismos del sistema. \\
\bottomrule
\end{tabularx}
\end{knowledgebox}

\teachervoice{00:09:53--00:11:14: Nathan define alignment de forma muy amplia: cabe en ella cualquier entrenamiento que acerque el modelo a los desires del usuario, mientras que IFT, SFT, RLHF y DPO son solo mecanismos distintos. Estas notas adoptan esa definición amplia, pero exigen precisar en cada caso los datos y la pérdida.}

\subsection{Chapter 0: la carrera inicial por reproducir ChatGPT}

Con la terminología establecida, el curso regresa a principios de 2023. En ese momento, la comunidad abierta no conocía la receta clave de los dialogue models y carecía de red-teaming, chat data y evaluaciones maduros. Numerosos equipos combinaron rápidamente base models, datos sintéticos y un fine-tuning ligero, lo que dio lugar a una carrera inicial de «un modelo por semana».

\lecturefigure{slide-027.jpg}{Chapter 0: el land grab por reproducir ChatGPT y sus problemas fundamentales}{official deck, p.~27}

\begin{warningbox}{No juzgar los primeros modelos con conocimiento retrospectivo}
El chat template, el system prompt, los datos de rechazo y la pairwise eval, que hoy parecen elementales, no tenían entonces una implementación unificada. El valor de los primeros resultados suele estar en que exponían problemas de interfaz y de datos, no en su puntuación final. El análisis histórico debe preguntar «qué desbloqueó», y no solo «si ya está obsoleto».
\end{warningbox}

\teachervoice{00:11:30--00:12:30: el ponente llama a este periodo land grab: todos publicaban modelos mientras apenas descubrían qué eran el dialogue, el red-teaming y una respuesta útil. Este contexto explica por qué los datos y la evaluación se convirtieron rápidamente en el campo de batalla principal.}

\subsection{Resumen de la sección}

Esta sección estableció la división en capas que necesita el análisis posterior: el base model determina la base de capacidades, IFT/SFT cambian la interfaz y la conducta de imitación, RLHF/DPO usan preferencias para reordenar las respuestas y alignment es un objetivo de sistema más amplio. El atlas del curso, por su parte, coloca en el mismo nivel los datos, la evaluación y los algoritmos, y nos recuerda que no debemos atribuir el avance de los modelos a un único optimizador.
